% ============================================================
%  APUNTES UNIVERSITARIOS — DERECHO PROCESAL CIVIL Y COMERCIAL
%  Autor: Isaac Edgar Camacho Ocampo
%  Generado con el Prompt Maestro Autónomo
% ============================================================

\documentclass[12pt,oneside]{book}

% ============================================================
% PAQUETES
% ============================================================
\usepackage[utf8]{inputenc}
\usepackage[T1]{fontenc}
\usepackage[spanish,es-nodecimaldot]{babel}

\usepackage[
    top=2.5cm,
    bottom=2.5cm,
    left=3cm,
    right=2.5cm,
    headheight=15pt
]{geometry}

\usepackage{amsmath}
\usepackage{amssymb}
\usepackage{mathtools}

\usepackage{graphicx}
\usepackage{float}
\usepackage{caption}
\usepackage{subcaption}

\usepackage{booktabs}
\usepackage{array}
\usepackage{longtable}
\usepackage{tabularx}

\usepackage{enumitem}

\usepackage{xcolor}
\usepackage[most]{tcolorbox}

\usepackage{fancyhdr}

\usepackage{ragged2e}

\usepackage[
    colorlinks=true,
    linkcolor=blue!50!black,
    citecolor=green!40!black,
    urlcolor=blue!60!black,
    pdftitle={Apuntes de Derecho Procesal Civil y Comercial},
    pdfauthor={Isaac Edgar Camacho Ocampo},
    pdfsubject={Apuntes universitarios},
    bookmarks=true
]{hyperref}

% ============================================================
% RUTAS DE IMÁGENES
% ============================================================
\graphicspath{
    {./img/}
    {./graficos/}
    {./figuras/}
}

% ============================================================
% CONFIGURACIÓN DE PÁRRAFOS Y LISTAS
% ============================================================
\setlength{\parindent}{0pt}
\setlength{\parskip}{0.5em}

\setlist[itemize]{topsep=0.3em, itemsep=0.2em}
\setlist[enumerate]{topsep=0.3em, itemsep=0.2em}

\clubpenalty=10000
\widowpenalty=10000

\setcounter{tocdepth}{2}

% ============================================================
% COMANDOS MATEMÁTICOS
% ============================================================
\newcommand{\abs}[1]{\left|#1\right|}
\newcommand{\norm}[1]{\left\lVert#1\right\rVert}
\newcommand{\vect}[1]{\mathbf{#1}}
\newcommand{\deriv}[2]{\frac{d#1}{d#2}}
\newcommand{\pderiv}[2]{\frac{\partial#1}{\partial#2}}

% ============================================================
% METADATOS
% ============================================================
\newcommand{\universidad}{Universidad de Buenos Aires (UBA)}
\newcommand{\facultad}{Facultad de Derecho}
\newcommand{\carrera}{}
\newcommand{\materia}{Derecho Procesal Civil y Comercial}
\newcommand{\titulo}{Apuntes de Derecho Procesal Civil y Comercial}
\newcommand{\autor}{Isaac Edgar Camacho Ocampo}
\newcommand{\anio}{2025}

% ============================================================
% CAJAS ACADÉMICAS
% ============================================================
\newtcolorbox{definition}[1]{
    enhanced, breakable,
    title={Definición: #1},
    fonttitle=\bfseries,
    colback=blue!3,
    colframe=blue!50!black,
    boxrule=0.8pt, arc=2mm
}

\newtcolorbox{important}{
    enhanced, breakable,
    title={Importante},
    fonttitle=\bfseries,
    colback=yellow!8,
    colframe=orange!70!black,
    boxrule=0.8pt, arc=2mm
}

\newtcolorbox{example}[1]{
    enhanced, breakable,
    title={Ejemplo: #1},
    fonttitle=\bfseries,
    colback=green!4,
    colframe=green!40!black,
    boxrule=0.8pt, arc=2mm
}

\newtcolorbox{formula}[1]{
    enhanced, breakable,
    title={Fórmula / Regla: #1},
    fonttitle=\bfseries,
    colback=purple!4,
    colframe=purple!50!black,
    boxrule=0.8pt, arc=2mm
}

\newtcolorbox{exercise}[1]{
    enhanced, breakable,
    title={Ejercicio: #1},
    fonttitle=\bfseries,
    colback=gray!5,
    colframe=gray!60!black,
    boxrule=0.8pt, arc=2mm
}

\newtcolorbox{note}{
    enhanced, breakable,
    title={Nota},
    fonttitle=\bfseries,
    colback=white,
    colframe=gray!50,
    boxrule=0.6pt, arc=2mm
}

% ============================================================
% ENCABEZADOS Y PIES
% ============================================================
\pagestyle{fancy}
\fancyhf{}
\fancyhead[L]{\nouppercase{\leftmark}}
\fancyhead[R]{\materia}
\fancyfoot[C]{\thepage}
\renewcommand{\headrulewidth}{0.4pt}
\renewcommand{\footrulewidth}{0pt}

\fancypagestyle{plain}{
    \fancyhf{}
    \fancyfoot[C]{\thepage}
    \renewcommand{\headrulewidth}{0pt}
}

% ============================================================
% INICIO DEL DOCUMENTO
% ============================================================
\begin{document}

% ============================================================
% PORTADA
% ============================================================
\begin{titlepage}
\thispagestyle{empty}

\begin{center}

\vspace*{1cm}

{\large \textsc{\universidad}}

\vspace{0.2cm}

{\large \textsc{\facultad}}

\vspace{3cm}

{\Huge \textbf{\materia}}

\vspace{1cm}

{\LARGE \titulo}

\vspace{0.8cm}

\textit{Comprender el proceso es el primer paso para ejercer el derecho con estrategia y responsabilidad.}

\vspace{1.2cm}

\rule{0.70\textwidth}{0.5pt}

\vspace{0.5cm}

{\large Apuntes de estudio}

\vspace{0.5cm}

\rule{0.70\textwidth}{0.5pt}

\vfill

{\large \autor}

\vspace{0.3cm}

{\large \anio}

\end{center}

\vfill

\begin{flushleft}
\textit{Este es un modesto aporte a los estudiantes de ``Derecho Procesal Civil y Comercial'' no pretende sustituir el material oficial de la materia.}
\end{flushleft}

\end{titlepage}

% ============================================================
% ÍNDICE
% ============================================================
\tableofcontents
\clearpage

% ============================================================
% CAPÍTULO 1: CADUCIDAD Y NEGLIGENCIA EN LA PRODUCCIÓN DE LA PRUEBA
% ============================================================
\chapter{Caducidad y Negligencia en la Producción de la Prueba}

\section{Contexto Procesal}

La caducidad y la negligencia en la producción de la prueba son dos institutos del derecho procesal que operan durante la etapa probatoria del proceso. Ambos se vinculan con el \textbf{principio dispositivo}: quien no cumple con una carga procesal puede sufrir consecuencias adversas a sus propios intereses.

\section{Caducidad de la Prueba}

\begin{definition}{Caducidad de la prueba}
La caducidad es la extinción de una carga procesal —entendida como facultad o acción— que se extingue por el mero transcurso de un \textbf{plazo estipulado en la ley}. Su elemento esencial es \textbf{objetivo}: el transcurso del tiempo fijado legalmente.
\end{definition}

La caducidad opera sobre los siguientes medios probatorios: prueba \textbf{pericial}, prueba \textbf{testimonial}, prueba \textbf{confesional} y prueba \textbf{documental}.

\subsection{Tipos de Caducidad}

\begin{table}[H]
\centering
\begin{tabularx}{\textwidth}{
    >{\bfseries\raggedright\arraybackslash}p{0.22\textwidth}
    >{\raggedright\arraybackslash}X
    >{\raggedright\arraybackslash}X
}
\toprule
\textbf{Tipo} & \textbf{Elemento} & \textbf{Alcance} \\
\midrule
Caducidad Objetiva  & Solo elemento objetivo (plazo) & Prueba pericial, testimonial, confesional y documental \\
\addlinespace
Caducidad Subjetiva & Elemento objetivo + conducta negligente & Tiene causa en un comportamiento omisivo del litigante \\
\bottomrule
\end{tabularx}
\caption{Tipos de caducidad en la producción de la prueba}
\end{table}

\section{Negligencia en la Producción de la Prueba}

\begin{definition}{Negligencia en la producción de la prueba}
La negligencia consiste en un \textbf{comportamiento omisivo} del litigante: falta de cuidado, falta de diligencia o falta de aplicación en la producción de la prueba. A diferencia de la caducidad, su elemento definitorio es \textbf{subjetivo}.
\end{definition}

\section{Diferencias entre Caducidad y Negligencia}

\begin{table}[H]
\centering
\begin{tabularx}{\textwidth}{
    >{\bfseries\raggedright\arraybackslash}p{0.22\textwidth}
    >{\raggedright\arraybackslash}X
    >{\raggedright\arraybackslash}X
}
\toprule
\textbf{Criterio} & \textbf{Caducidad} & \textbf{Negligencia} \\
\midrule
Elemento        & Objetivo (transcurso del plazo) & Subjetivo (conducta omisiva) \\
\addlinespace
Declaración     & De oficio o a pedido de parte & Solo a pedido de parte \\
\addlinespace
Operación       & Automática: el juez verifica el plazo & El juez aprecia la conducta del litigante \\
\addlinespace
Efecto          & Hacer caer la prueba de la parte contra quien se declara & Igual efecto \\
\bottomrule
\end{tabularx}
\caption{Cuadro comparativo: caducidad vs.\ negligencia}
\end{table}

\section{Consecuencias Procesales y Responsabilidad del Abogado}

\begin{important}
Las resoluciones relativas a la producción o denegación de prueba son \textbf{inapelables en forma inmediata}. Si el juez decreta una negligencia o una caducidad, la resolución queda firme. La responsabilidad por estas situaciones recae directamente sobre el abogado.
\end{important}

\subsection{Replanteo de Prueba en Segunda Instancia}

Cuando se dicta la sentencia definitiva, al apelarla, el apelante puede solicitar a la \textbf{cámara un replanteo de la prueba} que fue mal denegada o mal declarada en negligencia en primera instancia. Sin embargo, en esa instancia las posibilidades ya están considerablemente reducidas.

\subsection{El Oficio Reiteratorio}

En materia de prueba informativa, cuando una entidad —pública o privada— no contesta dentro del plazo legal, el abogado debe requerir el \textbf{oficio reiteratorio}. La inactividad frente a una respuesta que no llega expone a la parte a que la contraria acuse negligencia.

\begin{example}{Oficio sin respuesta y riesgo de negligencia}
Se libra un oficio a una entidad. La ley prevé un plazo de respuesta (por ejemplo, 20 días para entidades públicas). Si transcurre ese plazo sin respuesta y el abogado no solicita el reiteratorio, la parte contraria puede presentar un escrito acusando negligencia. Esa situación resulta difícil de sostener como defensa.
\end{example}

\subsection{Impacto sobre la Estrategia Probatoria}

Si la prueba afectada por caducidad o negligencia es la prueba central de la estrategia, la causa puede quedar seriamente comprometida. Por ello, el abogado debe:

\begin{itemize}
    \item Controlar permanentemente los plazos probatorios.
    \item Verificar el avance en la producción de la prueba propia.
    \item Monitorear la prueba de la parte contraria para poder, en su caso, solicitar la declaración de caducidad o negligencia.
\end{itemize}

% ============================================================
% CUADRO CORNELL — CAPÍTULO 1
% ============================================================
\section*{Cuadro Cornell: Caducidad y Negligencia}
\addcontentsline{toc}{section}{Cuadro Cornell: Caducidad y Negligencia}

\begin{table}[H]
\centering
\begin{tabularx}{\textwidth}{
    >{\raggedright\arraybackslash}p{0.30\textwidth}
    >{\raggedright\arraybackslash}X
}
\toprule
\textbf{Preguntas / palabras clave} & \textbf{Notas / respuestas} \\
\midrule

\textbf{¿Qué es la caducidad de la prueba?} &
Es la extinción de una facultad procesal por el transcurso del plazo legal. Opera automáticamente y puede declararse de oficio o a pedido de parte. \\
\addlinespace

\textbf{¿Qué es la negligencia en la prueba?} &
Es un comportamiento omisivo del litigante: falta de diligencia en impulsar la producción de la prueba. Tiene elemento subjetivo y solo se declara a pedido de parte. \\
\addlinespace

\textbf{¿Cuál es el efecto de ambas?} &
Hacer caer la prueba de la parte contra quien se declaran. La responsabilidad recae en el abogado. \\
\addlinespace

\textbf{¿Son apelables de inmediato?} &
No. Las resoluciones sobre producción o denegación de prueba son inapelables en forma inmediata. Solo pueden replantearse al apelar la sentencia definitiva. \\
\addlinespace

\textbf{¿Qué diferencia hay en la declaración?} &
La caducidad puede declararse de oficio o a pedido de parte. La negligencia solo puede declararse a pedido de parte. \\

\midrule

\multicolumn{2}{p{\textwidth}}{%
\textbf{Resumen:} La caducidad y la negligencia son los dos mecanismos por los que se puede perder la posibilidad de producir prueba. La primera opera automáticamente por el transcurso del plazo (elemento objetivo); la segunda requiere una conducta omisiva del litigante (elemento subjetivo) y solo se declara a pedido de parte. Ambas son inapelables de inmediato y comprometen directamente la responsabilidad del abogado. La única vía posterior es el replanteo de prueba ante la cámara al apelar la sentencia definitiva.
} \\

\bottomrule
\end{tabularx}
\end{table}

% ============================================================
% CAPÍTULO 2: HECHOS NUEVOS Y SISTEMA RECURSIVO
% ============================================================
\chapter{Hechos Nuevos, Recursos y Efecto Diferido}

\section{Concepto de Hechos Nuevos}

\begin{definition}{Hechos nuevos}
Son los \textbf{hechos jurídicamente relevantes} —ya sean simples hechos o actos jurídicos— que ocurren con posterioridad al intercambio de la demanda y la contestación de demanda, y que guardan relación directa con el objeto del litigio.
\end{definition}

\section{Ubicación Temporal y Plazo}

Los hechos nuevos se ubican en el siguiente tramo del proceso:

\begin{table}[H]
\centering
\begin{tabularx}{\textwidth}{
    >{\centering\arraybackslash}X
    >{\centering\arraybackslash}X
    >{\centering\arraybackslash\bfseries}X
    >{\centering\arraybackslash}X
}
\toprule
Demanda & Contestación de demanda & Hechos nuevos (hasta 5 días antes de la audiencia art.\ 360) & Audiencia art.\ 360 \\
\bottomrule
\end{tabularx}
\caption{Línea del tiempo: hechos nuevos en el proceso}
\end{table}

\begin{important}
Conforme al \textbf{artículo 363 del CPCCN}, la audiencia del artículo 360 se notifica mediante cédula. Los hechos nuevos pueden alegarse hasta el \textbf{quinto día} de notificada dicha audiencia. Adicionalmente, desde que se tomó conocimiento del hecho nuevo, el litigante cuenta con \textbf{cinco días} para denunciarlo.
\end{important}

\section{Traslado y Bilateralidad}

Según el artículo correspondiente, se le dará traslado a la otra parte sobre los hechos nuevos ``cuando se lo considere pertinente''. Esta formulación es ambigua. Sin embargo, desde el \textbf{principio de bilateralidad} y el \textbf{principio de defensa en juicio}, si no se corre traslado, la parte contraria no podrá refutar los hechos alegados. La falta de traslado implicaría, en la práctica, el rechazo de los hechos nuevos.

\begin{example}{Hecho nuevo en una causa de alimentos}
En una causa de alimentos, el demandado es despedido de su empleo con posterioridad a la presentación de la demanda. Ese despido constituye un hecho nuevo, ya que modifica la capacidad económica del demandado —base sobre la cual se determina la cuota— y guarda directa relación con el objeto del litigio. A la inversa, si el demandado mejorara sustancialmente su situación patrimonial, la parte actora también podría denunciarlo como hecho nuevo para fundamentar una cuota mayor.
\end{example}

\section{El Sistema Recursivo}

\subsection{El Agravio como Presupuesto del Recurso}

Para estar en condiciones de recurrir una resolución judicial, esta debe causar un \textbf{agravio}: un resultado diferente a lo pretendido o solicitado. Ante ese agravio, el primer paso es identificar el tipo de resolución de que se trata, porque de ello depende qué recurso procede.

\subsection{Pasos del Recurso de Apelación}

\begin{table}[H]
\centering
\begin{tabularx}{\textwidth}{
    >{\centering\arraybackslash}p{0.08\textwidth}
    >{\raggedright\arraybackslash}X
    >{\raggedright\arraybackslash}p{0.28\textwidth}
}
\toprule
\textbf{Paso} & \textbf{Acción} & \textbf{Plazo / Nota} \\
\midrule
1 & Interponer el recurso (escrito simple: ``vengo a apelar porque me causa agravio'') & 5 días (interlocutorias y definitivas) \\
\addlinespace
2 & El mismo juez de primera instancia concede o deniega el recurso & Análisis formal: tipo de resolución y plazo \\
\addlinespace
3A & Si concede: se fundan los agravios ante la cámara & Trámite inmediato o diferido \\
\addlinespace
3B & Si deniega: recurso de queja ante la cámara para que ordene conceder el recurso & Si no se va en queja: queda firme (cosa juzgada) \\
\bottomrule
\end{tabularx}
\caption{Pasos del recurso de apelación}
\end{table}

\subsection{El Recurso de Queja}

El recurso de queja se interpone ante la cámara cuando el juez de primera instancia denegó la apelación. Si la cámara hace lugar a la queja, el resultado es que el juez debe conceder la posibilidad de que la cámara trate el recurso de apelación. La queja, por sí misma, \textbf{no resuelve el fondo} de la cuestión; únicamente habilita el trámite del recurso principal. Si no se interpone queja y la resolución queda firme, se convierte en \textbf{cosa juzgada}.

\section{El Efecto Diferido}

\begin{definition}{Efecto diferido}
Es una modalidad de tramitación del recurso de apelación concedido, según la cual el recurso \textbf{no se eleva de inmediato a la cámara}, sino que queda en suspenso hasta que se dicte la sentencia definitiva. Solo se tramita si la parte que apeló pierde el juicio; si gana, el recurso diferido carece de interés práctico.
\end{definition}

\begin{note}
El término ``efecto'' en la expresión ``efecto diferido'' puede inducir a confusión. Conceptualmente, lo que se difiere no es el efecto del recurso, sino su \textbf{trámite} ante la cámara. La finalidad es evitar dilaciones innecesarias en el proceso: si el apelante obtiene una sentencia favorable, no necesitará que se trate ese recurso.
\end{note}

El efecto diferido se aplica, entre otros supuestos, cuando el juez rechaza un hecho nuevo y el litigante apela esa decisión.

\section{Tipos de Error Judicial en las Resoluciones Recurribles}

\begin{table}[H]
\centering
\begin{tabularx}{\textwidth}{
    >{\bfseries\raggedright\arraybackslash}p{0.30\textwidth}
    >{\raggedright\arraybackslash}X
}
\toprule
\textbf{Tipo de error} & \textbf{Descripción} \\
\midrule
Error \textit{in procedendo} & Violación de una norma procesal que afecta el derecho de defensa. Provoca la nulidad de la sentencia. \\
\addlinespace
Error \textit{in iudicando} (hechos) & El juez apreció incorrectamente los hechos: lo probado es diferente a lo que el juez afirmó en la sentencia. \\
\addlinespace
Error \textit{in iudicando} (derecho) & El juez aplicó una norma equivocada para resolver el caso. \\
\bottomrule
\end{tabularx}
\caption{Tipos de error judicial que fundamentan recursos}
\end{table}

% ============================================================
% CUADRO CORNELL — CAPÍTULO 2
% ============================================================
\section*{Cuadro Cornell: Hechos Nuevos, Recursos y Efecto Diferido}
\addcontentsline{toc}{section}{Cuadro Cornell: Hechos Nuevos, Recursos y Efecto Diferido}

\begin{table}[H]
\centering
\begin{tabularx}{\textwidth}{
    >{\raggedright\arraybackslash}p{0.30\textwidth}
    >{\raggedright\arraybackslash}X
}
\toprule
\textbf{Preguntas / palabras clave} & \textbf{Notas / respuestas} \\
\midrule

\textbf{¿Qué son los hechos nuevos?} &
Hechos jurídicamente relevantes ocurridos con posterioridad al intercambio de demanda y contestación, que guardan relación con el objeto del litigio. \\
\addlinespace

\textbf{¿Hasta cuándo se pueden alegar?} &
Hasta cinco días antes de la audiencia del art.\ 360, y dentro de los cinco días de tomado conocimiento del hecho nuevo (art.\ 363 CPCCN). \\
\addlinespace

\textbf{¿Por qué es obligatorio el traslado?} &
Por el principio de bilateralidad y defensa en juicio: la parte contraria debe poder refutar los hechos nuevos alegados. Sin traslado, el hecho nuevo debe considerarse rechazado. \\
\addlinespace

\textbf{¿Qué es el efecto diferido?} &
Modalidad de tramitación del recurso de apelación en la que este no sube de inmediato a la cámara, sino que queda pendiente hasta la sentencia definitiva. Solo se tramita si la parte apelante pierde el juicio. \\
\addlinespace

\textbf{¿Qué diferencia hay entre queja y apelación?} &
La apelación impugna el fondo de la resolución. La queja impugna la denegación del recurso de apelación y solo busca que la cámara ordene conceder ese recurso; no resuelve el fondo. \\
\addlinespace

\textbf{¿Cuáles son los errores que fundan un recurso?} &
Error \textit{in procedendo} (violación procesal), error \textit{in iudicando} en los hechos (mala apreciación), y error \textit{in iudicando} en el derecho (norma equivocada). \\

\midrule

\multicolumn{2}{p{\textwidth}}{%
\textbf{Resumen:} Los hechos nuevos son hechos sobrevinientes al intercambio de escritos que guardan relación con el litigio. Pueden alegarse hasta cinco días antes de la audiencia del art.\ 360 y requieren traslado a la contraria. Si el juez los rechaza, la apelación puede concederse con efecto diferido, lo que significa que no sube a la cámara de inmediato, sino al finalizar el proceso. El sistema recursivo se activa ante un agravio; el recurso de queja opera cuando se deniega la apelación, y únicamente habilita su trámite, sin resolver el fondo. Los errores que fundan recursos pueden ser procesales (\textit{in procedendo}) o sustanciales (\textit{in iudicando}).
} \\

\bottomrule
\end{tabularx}
\end{table}

% ============================================================
% CAPÍTULO 3: CUESTIÓN DE PURO DERECHO
% ============================================================
\chapter{Cuestión de Puro Derecho (Artículo 359 CPCCN)}

\section{Concepto y Base Legal}

\begin{definition}{Cuestión de puro derecho}
La cuestión de puro derecho es una \textbf{simplificación procesal} que opera cuando los hechos no son controvertidos entre las partes o cuando no existe prueba que producir. En esos casos, el debate se concentra exclusivamente en la \textbf{interpretación y aplicación de las normas}, sin que haya etapa probatoria.
\end{definition}

\begin{formula}{Artículo 359 CPCCN}
Habiendo sido contestada la demanda o la reconvención, y no quedando cuestiones previas que resolver, el juez tiene la facultad de resolver la cuestión de puro derecho.
\end{formula}

\section{Diferencia Central con el Proceso Ordinario}

En la cuestión de puro derecho \textbf{no se discute qué ocurrió} (los hechos), sino quién tiene mayor apoyo en la norma aplicable. La discusión es exclusivamente jurídica. No hay producción de prueba.

\section{Supuestos en que Procede}

Pueden identificarse los siguientes escenarios:

\begin{enumerate}
    \item El demandado admite que los hechos ocurrieron tal como los describe la parte actora.
    \item Las partes no tienen pruebas que ofrecer y aceptan continuar sin etapa probatoria.
    \item En la \textbf{audiencia preliminar del artículo 360}, el juez evalúa la situación y declara que los hechos están claros, configurando una cuestión que debe resolverse de puro derecho.
\end{enumerate}

\section{Procedimiento}

Una vez declarada la cuestión de puro derecho:

\begin{enumerate}
    \item El juez concede un plazo de \textbf{5 días} a la parte que inició el proceso para que se expida.
    \item Se corre traslado a la parte contraria por igual plazo de \textbf{5 días}.
    \item Las partes presentan sus \textbf{alegatos}.
    \item Si la resolución produce agravio, es \textbf{apelable}.
    \item De no mediar apelación, el juez dicta sentencia.
\end{enumerate}

\section{Aplicación en el Juicio Ejecutivo}

En el \textbf{juicio ejecutivo}, el título ejecutivo que porta el accionante otorga presunción de legitimidad. El juez da traslado al demandado por \textbf{10 días} para que oponga excepciones. Si no las hay o son rechazadas, se dicta la \textbf{sentencia de remate}. Este trámite también configura una cuestión de puro derecho en tanto no se discuten los hechos subyacentes sino la validez del título.

% ============================================================
% CAPÍTULO 4: ACCIÓN MERAMENTE DECLARATIVA E INCONSTITUCIONALIDAD
% ============================================================
\chapter{Acción Meramente Declarativa e Inconstitucionalidad}

\section{Concepto de Acción Meramente Declarativa}

\begin{definition}{Acción meramente declarativa}
Es la acción judicial mediante la cual se solicita al órgano jurisdiccional que \textbf{aclare o declare una situación jurídica}, sin que medie necesariamente una condena de cumplimiento inmediato. Sirve para eliminar una incertidumbre sobre el alcance, la existencia o la modalidad de una relación jurídica o de la validez de una norma.
\end{definition}

La acción meramente declarativa es la herramienta idónea para impugnar la \textbf{constitucionalidad de una norma} cuando existe un caso concreto, sin necesidad de esperar un daño consumado.

\section{Requisito del Caso Concreto}

Para plantear la inconstitucionalidad de una norma mediante acción meramente declarativa, debe existir un \textbf{caso concreto}: una parte afectada de manera directa, actual y específica por la aplicación de la norma cuestionada. No es posible plantear la inconstitucionalidad en abstracto.

\section{Caso Práctico: El Consejo Federal Pesquero}

\subsection{Contexto}

El \textbf{Consejo Federal Pesquero} es un organismo creado por ley en el año 1998. Está integrado por representantes de las provincias con costa marina (Buenos Aires, Río Negro, Chubut, Santa Cruz y Tierra del Fuego) y por representantes del Ministerio de Economía, Agricultura y Pesca. Entre sus atribuciones, la ley le otorga la facultad de determinar anualmente qué especies se pescarán y de asignar cuotas de pesca a cada empresa.

\subsection{El Problema Jurídico: Ausencia de Licitación}

Una empresa pesquera es excluida del otorgamiento de la cuota de merluza sin que mediara un procedimiento licitatorio. En el marco de las contrataciones públicas, la adquisición o concesión de bienes del Estado requiere un procedimiento de \textbf{licitación pública}: un proceso transparente por el cual se convoca a distintos oferentes y se adjudica al que mejor satisfaga los requisitos establecidos, en términos de calidad y precio.

En el caso, el Consejo Federal Pesquero asignó las cuotas de manera discrecional, sin licitación, lo que plantea una posible violación al principio de \textbf{igualdad de condiciones} y a los principios que rigen el uso de bienes públicos.

\subsection{Estrategia Jurídica}

La estrategia adoptada fue iniciar una \textbf{acción meramente declarativa de inconstitucionalidad}, solicitando al juez que declare la inconstitucionalidad de la ley que crea el Consejo Federal Pesquero y de su reglamentación, en cuanto otorgan atribuciones incompatibles con la Constitución Nacional.

\begin{important}
El objeto de la acción no es obtener la cuota de pesca directamente, sino lograr que se declare la inconstitucionalidad del sistema de asignación discrecional, de modo que el Estado quede obligado a establecer un mecanismo transparente —licitatorio— en el que todos los interesados puedan participar en igualdad de condiciones.
\end{important}

\subsection{Planteo de Inconstitucionalidad}

El planteo ante el juez se estructura en los siguientes términos: la ley dictada por el Congreso es inconstitucional porque genera un organismo al que atribuye facultades de asignación discrecional de recursos públicos; y la reglamentación del Ejecutivo agrava esa inconstitucionalidad. Quien ejerce una actividad lícita tiene derecho a acceder a los bienes públicos en igualdad de condiciones con los demás actores del sector.

\subsection{Demandados}

Se demanda:
\begin{itemize}
    \item Al \textbf{Estado Nacional}, como responsable del acto legislativo (la ley del Congreso).
    \item Al \textbf{Poder Ejecutivo}, por haber reglamentado la ley de un modo que agrava la inconstitucionalidad.
\end{itemize}

La estrategia defensiva del Estado fue cuestionar su legitimación pasiva, sosteniendo que el legitimado pasivo era el Consejo Federal Pesquero y no el Estado Nacional directamente. La solución del juez fue incorporar al Consejo Federal Pesquero como \textbf{codemandado}.

\section{La Cuestión de Puro Derecho en este Caso}

\begin{exercise}{Tarea de análisis}
Una vez contestada la demanda por el Consejo Federal Pesquero, ¿habrá o no habrá audiencia del artículo 360? ¿Y por qué? Justificar desde el tipo de controversia planteada.
\end{exercise}

En este caso, la controversia es de naturaleza \textbf{exclusivamente normativa}: la discusión se centra en si la ley y su reglamentación son o no inconstitucionales. No hay hechos empíricos que probar. No existen hechos controvertidos entre las partes respecto a cómo opera el sistema de cuotas: lo que se discute es si ese sistema es compatible con la Constitución Nacional. Por lo tanto, la causa se convierte en una \textbf{cuestión de puro derecho}, sin etapa probatoria relevante.

Adicionalmente, al tratarse de una acción de inconstitucionalidad, no es posible proponer una conciliación en la audiencia preliminar, ya que si la controversia versa sobre la validez constitucional de una norma, no hay margen de acuerdo entre las partes sobre ese punto.

% ============================================================
% CUADRO CORNELL — CAPÍTULOS 3 Y 4
% ============================================================
\section*{Cuadro Cornell: Cuestión de Puro Derecho y Acción Meramente Declarativa}
\addcontentsline{toc}{section}{Cuadro Cornell: Cuestión de Puro Derecho y Acción Meramente Declarativa}

\begin{table}[H]
\centering
\begin{tabularx}{\textwidth}{
    >{\raggedright\arraybackslash}p{0.30\textwidth}
    >{\raggedright\arraybackslash}X
}
\toprule
\textbf{Preguntas / palabras clave} & \textbf{Notas / respuestas} \\
\midrule

\textbf{¿Cuándo procede la cuestión de puro derecho?} &
Cuando los hechos no son controvertidos (el demandado los admite), cuando las partes no tienen prueba que ofrecer, o cuando el juez así lo declara en la audiencia del art.\ 360. La discusión es exclusivamente jurídica. \\
\addlinespace

\textbf{¿Cómo se tramita?} &
El juez concede cinco días a la actora; luego traslado de cinco días a la contraria; las partes presentan alegatos; y se dicta sentencia. La resolución es apelable. \\
\addlinespace

\textbf{¿Qué diferencia hay con el juicio ordinario?} &
No se discute qué ocurrió (los hechos), sino quién tiene mayor apoyo en la norma. No hay producción de prueba. \\
\addlinespace

\textbf{¿Qué busca la acción meramente declarativa?} &
Aclarar o declarar una situación jurídica incierta sin condena de cumplimiento inmediato. Es la vía para plantear la inconstitucionalidad de una norma en un caso concreto. \\
\addlinespace

\textbf{¿Requiere caso concreto?} &
Sí. No puede plantearse la inconstitucionalidad en abstracto. Debe haber una parte directamente afectada por la norma cuestionada. \\
\addlinespace

\textbf{¿Por qué el caso pesquero es de puro derecho?} &
Porque la controversia es exclusivamente normativa: si la ley y su reglamentación que regulan las cuotas de pesca son inconstitucionales. No hay hechos empíricos que probar ni espacio para conciliación. \\

\midrule

\multicolumn{2}{p{\textwidth}}{%
\textbf{Resumen:} La cuestión de puro derecho (art.\ 359 CPCCN) opera cuando los hechos no son controvertidos y la discusión es exclusivamente jurídica: no hay prueba, solo interpretación normativa. La acción meramente declarativa es el instrumento para obtener una declaración de certeza sobre una situación jurídica incierta —incluida la inconstitucionalidad de una norma— cuando existe un caso concreto. El caso del Consejo Federal Pesquero ilustra ambas figuras: se impugna un sistema de asignación discrecional de cuotas de pesca sin licitación, y al no haber hechos que probar, la causa deviene en cuestión de puro derecho.
} \\

\bottomrule
\end{tabularx}
\end{table}

% ============================================================
% CAPÍTULO 5: ESTRATEGIA PROBATORIA EN RESPONSABILIDAD CIVIL
% ============================================================
\chapter{Estrategia Probatoria en Responsabilidad Civil}

\section{El Caso Práctico: Accidente en Embarcación}

El caso analizado involucraba un accidente producido en una embarcación en el que falleció una persona. La responsabilidad del \textbf{patrón} fue caracterizada como subjetiva, mientras que la responsabilidad de la \textbf{empresa de catamaranes} (como tercero contratante) fue calificada como objetiva.

\section{El Eje Central: Probar el Hecho vs.\ Probar los Daños}

\begin{important}
Para reclamar daños en sede judicial, deben probarse \textbf{dos cosas distintas}:
\begin{enumerate}
    \item El \textbf{hecho causante} del daño.
    \item Los \textbf{daños concretos} sufridos como consecuencia de ese hecho.
\end{enumerate}
La prueba del hecho causante no es suficiente para acreditar los daños. Sin demostrar el ``antes'' y el ``después'', no hay daño acreditable.
\end{important}

\begin{definition}{Daño como cambio}
El daño consiste en la diferencia entre la situación de la persona \textbf{antes} del hecho y la situación \textbf{después} de él. Esa brecha —física, psicológica, patrimonial o social— es lo que debe ser acreditado con precisión.
\end{definition}

\section{Medios de Prueba del Hecho Causante}

\subsection{Prueba Testimonial}

Se propusieron como testigos:
\begin{itemize}
    \item Personas que se encontraban en el catamarán al momento del accidente, que habrían escuchado bocinas y pedidos de auxilio.
    \item Un testigo clave: una persona que se giró en el momento del impacto y pudo observar directamente lo ocurrido.
\end{itemize}

\subsection{Prueba Informativa}

Se propuso librar oficios a:
\begin{itemize}
    \item La municipalidad, para obtener registros de cámaras de seguridad.
    \item Gendarmería Nacional o Prefectura Naval Argentina, por las cámaras de biovigilancia y registros sobre el hecho.
    \item El organismo de emergencias médicas que atendió a la víctima (SAME u organismo equivalente en la jurisdicción de Tigre), por sus registros de atención.
    \item La aseguradora correspondiente, para obtener la póliza o confirmar su existencia.
\end{itemize}

\subsection{Prueba Documental}

Se propuso ofrecer:
\begin{itemize}
    \item La libreta de casamiento (para acreditar el vínculo matrimonial con la víctima).
    \item El acta constitutiva del estudio jurídico del fallecido y el libro contable (para acreditar sus ingresos).
    \item El canal de habilitación del catamarán ante la Prefectura Naval.
    \item El carnet náutico del patrón. Si no se contara con él, podría obtenerse mediante oficio al club náutico correspondiente o a Prefectura.
    \item Comprobantes de gastos médicos y de clínica.
    \item El expediente penal o la instrucción penal preparatoria (IPP) iniciada con motivo del fallecimiento.
\end{itemize}

\subsection{Prueba Pericial}

Se mencionó un perito mecánico especialista en embarcaciones para analizar las circunstancias técnicas del choque. Se observó que podría resultar sobreabundante dado que otros medios de prueba ya cubrían ese aspecto.

\section{Medios de Prueba de los Daños}

\subsection{Daño Moral}

Para acreditar el daño moral sufrido por el cónyuge sobreviviente, puede ofrecerse prueba testimonial de personas conocidas que puedan dar cuenta de:
\begin{itemize}
    \item La vida social y afectiva de la pareja antes del accidente.
    \item Los cambios observados en la persona afectada con posterioridad al hecho.
\end{itemize}

\begin{example}{Testimonial para daño moral}
Testigos que frecuentaban el club con la pareja pueden declarar: ``Venían al club todos los días como pareja; ahora este señor ya no viene más''. Ese tipo de testimonio permite mostrar concretamente lo que la persona ha perdido como consecuencia del accidente.
\end{example}

\subsection{Daño Psicológico}

El accidente puede haber desencadenado una depresión, un trastorno de estrés postraumático u otras afectaciones psicológicas, dependiendo de la evolución de cada persona.

Para reclamar daño psicológico resarcible:
\begin{itemize}
    \item Debe ofrecerse una \textbf{pericia psicológica}.
    \item Deben formularse los \textbf{puntos de pericia} adecuados para acreditar el daño psicológico concreto y su relación causal con el hecho.
\end{itemize}

\subsection{Pérdida de Chance y Lucro Cesante}

Para reclamar la pérdida de chance y el lucro cesante derivados del fallecimiento del cónyuge:
\begin{itemize}
    \item Ofrecer las \textbf{facturaciones} del estudio jurídico (en este caso, más que recibos de sueldo, dada la actividad profesional independiente).
    \item Ofrecer prueba testimonial de compañeros de trabajo o socios que puedan acreditar la actividad profesional del fallecido, su productividad y sus ingresos habituales.
\end{itemize}

\begin{example}{Testimonial para pérdida de chance}
Testimonios de socios o compañeros de trabajo que acrediten: ``Eran socios hace diez años; trabajaban juntos y participaban de la vida social. Ahora, a raíz del hecho, el sobreviviente ya no rema, no navega, le tiene miedo al agua''. Esa clase de testimonio aporta elementos para cuantificar la pérdida.
\end{example}

% ============================================================
% CUADRO CORNELL — CAPÍTULO 5
% ============================================================
\section*{Cuadro Cornell: Estrategia Probatoria en Responsabilidad Civil}
\addcontentsline{toc}{section}{Cuadro Cornell: Estrategia Probatoria en Responsabilidad Civil}

\begin{table}[H]
\centering
\begin{tabularx}{\textwidth}{
    >{\raggedright\arraybackslash}p{0.30\textwidth}
    >{\raggedright\arraybackslash}X
}
\toprule
\textbf{Preguntas / palabras clave} & \textbf{Notas / respuestas} \\
\midrule

\textbf{¿Qué dos cosas deben probarse para reclamar daños?} &
El hecho causante y los daños concretos sufridos. Probar el hecho no es suficiente; debe acreditarse el cambio en la situación de la persona (``antes'' y ``después''). \\
\addlinespace

\textbf{¿Cómo se acredita el daño moral?} &
Con prueba testimonial de personas que puedan dar cuenta de la vida de la persona antes del hecho y de las pérdidas concretas que sufrió como consecuencia. \\
\addlinespace

\textbf{¿Cómo se acredita el daño psicológico?} &
Con una pericia psicológica cuyos puntos de pericia estén formulados de modo que acrediten el daño resarcible y su relación causal con el hecho. \\
\addlinespace

\textbf{¿Qué prueba se necesita para pérdida de chance?} &
Las facturaciones o registros de ingresos del fallecido, más prueba testimonial de compañeros de trabajo o socios que acrediten la actividad y los ingresos habituales. \\
\addlinespace

\textbf{¿Qué es el daño como diferencia?} &
El daño es la brecha entre la situación de la persona antes del hecho y su situación posterior. Sin demostrar esa diferencia, no hay daño acreditable en juicio. \\
\addlinespace

\textbf{¿Qué tipos de responsabilidad se distinguen en el caso?} &
Responsabilidad subjetiva del patrón (debe probarse la culpa) y responsabilidad objetiva de la empresa de catamaranes (no requiere demostrar culpa). \\

\midrule

\multicolumn{2}{p{\textwidth}}{%
\textbf{Resumen:} La estrategia probatoria en responsabilidad civil requiere distinguir dos planos: la prueba del hecho causante (testimonial, informativa, documental, pericial) y la prueba de los daños concretos (moral, psicológico, patrimonial). Para el daño moral, la testimonial es central; para el daño psicológico, se requiere pericia psicológica con puntos adecuados; para la pérdida de chance y el lucro cesante, corresponde acreditar los ingresos y la actividad del fallecido. El concepto de daño como diferencia entre el ``antes'' y el ``después'' es el eje de toda la estrategia.
} \\

\bottomrule
\end{tabularx}
\end{table}

% ============================================================
% CUADRO CORNELL INTEGRADOR FINAL
% ============================================================
\chapter*{Cuadro Cornell Integrador}
\addcontentsline{toc}{chapter}{Cuadro Cornell Integrador}

\begin{table}[H]
\centering
\begin{tabularx}{\textwidth}{
    >{\raggedright\arraybackslash}p{0.30\textwidth}
    >{\raggedright\arraybackslash}X
}
\toprule
\textbf{Preguntas / palabras clave} & \textbf{Notas / respuestas} \\
\midrule

\textbf{¿Qué es la caducidad de la prueba?} &
Extinción de una facultad procesal por el transcurso del plazo legal. Opera automáticamente; puede declararse de oficio o a pedido de parte. Afecta prueba pericial, testimonial, confesional y documental. \\
\addlinespace

\textbf{¿Qué es la negligencia en la prueba?} &
Comportamiento omisivo del litigante en la producción de la prueba. Elemento subjetivo. Solo se declara a pedido de parte. \\
\addlinespace

\textbf{¿Cuál es el efecto de ambas?} &
Hacer caer la prueba de la parte contra quien se declaran. La responsabilidad recae directamente sobre el abogado. Las resoluciones son inapelables en forma inmediata. \\
\addlinespace

\textbf{¿Qué son los hechos nuevos?} &
Hechos jurídicamente relevantes ocurridos con posterioridad al intercambio de demanda y contestación, relacionados con el objeto del litigio. Plazo: hasta cinco días antes de la audiencia del art.\ 360 (art.\ 363 CPCCN). \\
\addlinespace

\textbf{¿Qué es el efecto diferido?} &
Modalidad por la que el recurso concedido no sube de inmediato a la cámara, sino que queda pendiente hasta la sentencia definitiva. Solo se tramita si la parte apelante pierde. Evita dilaciones. \\
\addlinespace

\textbf{¿Cuándo procede la cuestión de puro derecho?} &
Cuando los hechos no son controvertidos o no hay prueba que producir. La controversia es exclusivamente jurídica: interpretación y aplicación de normas. El juez puede declararlo en la audiencia del art.\ 360. \\
\addlinespace

\textbf{¿Qué es la acción meramente declarativa?} &
Acción judicial que busca declarar una situación jurídica incierta sin condena inmediata. Es la vía idónea para impugnar la constitucionalidad de una norma en un caso concreto. \\
\addlinespace

\textbf{¿Cómo se acredita un daño?} &
Deben probarse dos componentes: el hecho causante y el cambio producido en la persona o el patrimonio (diferencia entre el ``antes'' y el ``después''). Sin acreditar esa diferencia, no hay daño resarcible. \\
\addlinespace

\textbf{¿Cómo funciona el sistema recursivo?} &
Ante un agravio: se interpone recurso en cinco días; el juez concede o deniega; si concede, se fundan agravios ante la cámara; si deniega, queja ante la cámara. Sin queja: cosa juzgada. \\
\addlinespace

\textbf{¿Por qué el caso pesquero es de puro derecho?} &
La controversia es normativa: si la ley y su reglamentación que regulan las cuotas de pesca son inconstitucionales. No hay hechos empíricos que probar. No cabe conciliación. \\

\midrule

\multicolumn{2}{p{\textwidth}}{%
\textbf{Resumen integrador:} La clase articuló cinco ejes del proceso civil ordinario. \textbf{Primero:} la caducidad y la negligencia son los dos mecanismos por los cuales un abogado puede perder la prueba de su cliente; la primera opera automáticamente por el transcurso del tiempo (elemento objetivo); la segunda, por una conducta omisiva (elemento subjetivo). Ambas son inapelables en forma inmediata y comprometen la responsabilidad del letrado. \textbf{Segundo:} los hechos nuevos permiten incorporar al proceso hechos sobrevinientes relevantes, sujetos a plazos estrictos y al principio de bilateralidad; los recursos contra su rechazo pueden tramitarse con efecto diferido, quedando pendientes hasta la sentencia definitiva para evitar dilaciones. \textbf{Tercero:} la cuestión de puro derecho (art.\ 359 CPCCN) permite resolver el litigio sin etapa probatoria cuando los hechos no son controvertidos, concentrando el debate en la interpretación normativa. \textbf{Cuarto:} la acción meramente declarativa es la herramienta para impugnar normas inconstitucionales ante un caso concreto; el caso del Consejo Federal Pesquero ilustra cómo un sistema de asignación discrecional de cuotas sin licitación puede impugnarse por esta vía, configurando además una cuestión de puro derecho. \textbf{Quinto:} la estrategia probatoria en responsabilidad civil exige distinguir siempre la prueba del hecho causante de la prueba de los daños concretos; el daño es la diferencia entre el ``antes'' y el ``después'', y esa brecha debe acreditarse con precisión mediante testimonial, pericial y documental.
} \\

\bottomrule
\end{tabularx}
\end{table}

\end{document}
